\lecture{2}

\section{Самый простой случай задачи оптимизации}
\begin{problem}[Задача оптимизации]
    Будем искать \(\min_{x \in \R^d} f(x)\), т.е. безусловный минимум
\end{problem}

\begin{definition}
    Точка \(x^*\) называется локальным минимумом функции \(f\) на \(\R^d\) (локальным решением задачи минимизации), если существуует такое \(r > 0\), что 
    \[\forall y \in B_r(x^*): f(x) \le f(y)\]
\end{definition}

\begin{definition}
    Точка \(x^*\) называется глобальным минимумом функции \(f\) на \(\R^d\) (глобальным решением задачи минимизации), если существуует такое \(r > 0\), что 
    \[\forall y \in \R^d: f(x) \le f(y)\]
\end{definition}

\begin{theorem}
    Пусть \(f: \R^d \ra \R\) --- дифференцируема, \(x^*\) --- ее локальный минимум. Тогда \(\nabla f(x^*) = 0\).
\end{theorem}
\begin{proof}
    См. Матан 3.
\end{proof}

\begin{definition}
    Пусть \(f\) --- непрерывно дифференцируемая функция на \(\R^d\). Будем говорить, что она выпуклая, если \(\forall x \)
\end{definition}
